\documentclass{article}
\usepackage[utf8]{inputenc}

\title{Template Tesis ITAM}
\author{Lorena Mejía Domenzain}
\date{February 2019}

\begin{document}

\maketitle

%\section{Introduction}

Template de LaTeX para una tesis del ITAM

En este template de LaTeX hay dos modos de hacer la tesis.
\begin{itemize}
    \item Por Capítulos: 
    
    Los capítulos se escriben cada uno en un archivo dentro del folder \emph{Chapters} y se compila el archivo \emph{porCapitulos.tex}.
    
    En \emph{porCapitulos.tex} se modifica la portada, el resumen y el aviso de divulgación.
    
    \item Todo junto:
    
    Escribes toda tu tesis en un mismo archivo \emph{todoJunto.tex}. Ahí esta la portada, el aviso de divulgación y los capítulos. Todo es un gran chorizo, como si fuera un documento de Word.
    
    \emph{todoJunto.tex} es el único archivo que hay que cuidar y compilar.
    
\end{itemize}

En ambos casos, las referencias van en el archivo \emph{Bibliography.bib}.

\vspace{3pt}

Ninguno de los dos tipos es mejor o peor. Yo utilice uno para mi tesis de licenciatura y otro para mi tesis de maestría. (Por eso mismo no son exactamente iguales en contenido).

\end{document}
